\section{Computational Evidence}\label{sec:computational}
\subsection{Comparisons, measurement, and audit scopes}
All instances are synthetic rational service agreements. The binary comparison family uses discount and payment coefficients one, two public regimes per noninitial date, independent seeded quadratic rewards, and continuation caps above a known feasible constant-tier policy. Thus an exact root promise is available without using any optimizer. The horizon ranges from four to ten dates for the unfolded comparisons, with two seeds at each horizon. A separate public-graph grid comparison extends to twelve dates. Every input and implemented policy is retained.

We implement four explicitly defined methods. The quotient compiler merges coefficient events on the public graph. The independent tree method unfolds the same input, changes variables as in Proposition~\ref{prop:laminar}, and performs occurrence-indexed exact marginal reduction without public-vertex memoization. It uses a separate interpolation implementation, not the quotient code. The generic numerical comparison applies SciPy's sparse trust-region interior-point method to the unfolded convex quadratic problem \citep{Virtanen2020}. The primal promise method performs backward dynamic programming directly on the public graph, restricting successor promises to a grid while allowing each current action to be the continuous residual. None of these baselines receives the quotient optimum as an input.

Each case--method pair runs in a fresh sequential worker. Three repetitions are retained with their minimum, median, and maximum. Numerical-library imports are outside timed method calls and are performed in all workers. Peak memory is an absolute process high-water mark, including the common runtime and libraries; it does not estimate allocator-level incremental memory. The comparison suite repeats construction, root inversion, certificate generation, and the inherited whole-response audit separately. In the scale suite, construction, inversion, and machine minimization are repeated three times; certificate generation and an independent single-query policy audit are each timed once and explicitly identified as diagnostics. The latter checks conditional caps, exact state flow, stationarity, complementarity, root payment, and reward without checking every affine response segment. The companion specifies all tolerances, generators, and measurement boundaries.

\subsection{What the baselines establish after the R31 reframing}
The exact tree comparison is a correctness cross-check, not an algorithmic novelty claim. All eight reported tree optima equal the quotient optima as exact rational numbers. The numerical convex solves terminate successfully; their largest absolute value difference is $3.64\times10^{-7}$ and their largest primal constraint residual is below $2.63\times10^{-14}$. These residuals are not exact optimality certificates.

For the implemented occurrence-indexed code and these inputs, the tree-to-quotient timing ratio ranges from approximately $1.9$ to $40.7$. R31 does not use those ratios as evidence against classical laminar algorithms. A normalized memoized occurrence reduction has the same recurrence as the quotient compiler. The computational conclusion from this comparison is exact agreement between two independently implemented representations.

The public-graph promise method is the more informative structural comparator because it already exploits recombination. Its policies satisfy the exact root promise and every continuation cap under rational replay. At resolution 80, the largest absolute loss among the six reported horizon--seed pairs is below $9.77\times10^{-5}$; one instance is already exact at resolution 40. The comparison isolates what the exact dual path contributes beyond graph compression: continuous parametric closure, exact root inversion, and a finite exact execution machine without selecting a promise mesh.

The modern parametric literature discussed in Section~\ref{sec:model} is handled theorem by theorem rather than through a mismatched software race. \citet{KlimmWarode2021}, for example, uses a parametric flow input and different parameterization. R31 therefore does not infer practical superiority from a generic QP timing or from an unfolded-tree count. Its computational questions are instead whether the compact compiler reproduces the exact laminar optimum, whether the whole response representation exhibits the proved output-size regimes, and whether one compiled path supports cheap repeated root queries.

\subsection{Representation tightness, graph scale, and rational precision}
Theorem~\ref{thm:tight-response} supplies an analytic worst-case family: an $n$-vertex chain has exactly $2n$ global knots and $n^2+2n$ stored affine response segments. This lower bound is independent of implementation and closes the one-sided complexity issue raised by the referee.

The separately retained dense-knot experiment provides an empirical counterpart. With slack continuation caps and distinct local clipping thresholds, a 511-vertex instance stores 197,369 affine segments, compared with 3,161 at 63 vertices. Coefficient lengths remain small and execution needs one writable symbol. The experiment is not used as the proof of quadratic storage; it illustrates the representation regime already established analytically.

A complementary late-renewal family retains heterogeneous binding caps at the penultimate date and recombines afterward. It reaches 2,045 public vertices and 8,164 edges. It has ten global response knots, 22,419 stored segments, and largest coefficient length 2,213 bits. The optimal protocol requires four writable symbols. Construction takes a median 2.62 seconds, while separate certificate generation and single-query audit measurements take 0.87 and 1.19 seconds. This is a long rational-horizon example with sparse cap activation, not a worst-case density test.

The broadly capped coupled family has heterogeneous payment and reward coefficients and rational transitions. Its maximum coefficient length rises from 198 to 388 to 768 bits as the horizon doubles from 16 to 32 to 64 dates; a tighter-cap 128-date instance reaches 1,501 bits. These observations are consistent with, but do not prove, the polynomial bit-height theorem. The theorem itself is algebraic and includes sorting, exact comparisons, and fraction reduction.

The response representation, the reached behavior machine, and the certificate are measured separately. A small writable alphabet does not imply a small read-only response table, and a sparse knot universe does not imply short aggregate certificate numerators. This is why R31 reports $K_*$, $C_*$, response segments, coefficient precision, and certificate diagnostics as distinct resources.

\subsection{Repeated root queries and the value of compiling the path}
A compiled root response supports repeated exact promise inversions without rebuilding all response tables. The retained experiment performs 101 evenly spaced root queries for each late-renewal size. These timings include inversion only; emitting and certifying a new policy for each promise is an additional operation. Corollary~\ref{cor:root-query} gives the corresponding exact complexity statement: after preprocessing, root inversion needs logarithmically many exact comparisons in the number of stored knots and a constant number of rational operations.

This repeated-query experiment is aligned with the residual parametric claim. The paper does not say that every application needs the entire path. It establishes when the complete path is compact enough to compile exactly and when that preprocessing can be amortized across root commitments or contractual design queries.

\subsection{Verification and the retained evidence boundary}
All 35 inherited exact cases reproduce their complete response curves and values, with 147,157 assertion groups in the inherited independent verifier. Forty additional small graph instances agree exactly with the independently implemented tree reduction. Two hundred and two comparisons of explicitly expanded continuation plans check the behavioral equivalence; a separate case reduces two distinct incoming prices to a single necessary symbol. Exhaustive enumeration checks the ordered frontier against 2,308 arbitrary set partitions, including unequal probabilities and nonquadratic costs. Boundary tests cover 150 feasible corridor anchors, 96 tied-normal occurrences, and 1,185 exact supporting inequalities. The test suite also rejects deliberately invalid inputs and altered certificates. These are implementation and finite-instance checks, not independent mathematical proofs of universal claims.

The earlier deployment study is neither removed nor used as favorable algorithmic evidence. At the strict tolerance, 83 of 84 cache queries refreshed and no cache speed advantage was measured. Switching paths, comparator-adjusted rents, continuous-state envelopes, correlated certificates, and governance results retain their original statements, assumptions, proofs, and numerical records in the supplementary theory volumes. The R31 response-complexity theorem and the earlier cache study address different algorithms and different questions.